\section*{Глава 1. Анализ предметной области}
\addcontentsline{toc}{section}{Глава 1. Анализ предметной области}
\stepcounter{section}

\subsection{Обзор существующих веб-сервисов для управления гардеробом}

На современном рынке цифровых решений для управления личным гардеробом выделяются приложения, предназначенные для каталогизации одежды, формирования образов и персонализированных рекомендаций пользователю. Такие сервисы позволяют систематизировать элементы гардероба и упрощают процесс подбора одежды в повседневной жизни.

Сравнительный анализ популярных решений представлен в таблице~\ref{tab:wardrobe_services}.

\begin{table}[ht!]
	\centering
	\caption{Сравнительный анализ веб-сервисов для управления гардеробом}
	\label{tab:wardrobe_services}
	\begin{tabular}{|p{0.22\textwidth}|p{0.48\textwidth}|p{0.22\textwidth}|}
		\hline
		\textbf{Сервис} & \textbf{Функциональные возможности} & \textbf{Модель доступа} \\ \hline
		
		Stylebook &
		Ручная каталогизация одежды, создание образов и планирование гардероба &
		Платное приложение \\ \hline
		
		Smart Closet &
		Ведение цифрового гардероба, создание образов и базовая автоматизация подбора одежды &
		Условно-бесплатная модель \\ \hline
		
		Cladwell &
		Рекомендации образов на основе предпочтений пользователя, сезона и погодных условий &
		Условно-бесплатная модель \\ \hline
		
		Wardrobe AI &
		Автоматический подбор одежды и формирование комплектов с использованием алгоритмов анализа данных &
		Условно-бесплатная модель \\ \hline
		
		Pureple &
		Визуальное управление гардеробом и создание капсульных коллекций &
		Бесплатное приложение \\ \hline
		
	\end{tabular}
\end{table}

Анализ показывает, что большинство существующих решений ориентировано на мобильные приложения и ограничено в возможностях расширенной персонализации. Также часто отсутствует гибкость в настройке логики формирования образов и интеграции пользовательских алгоритмов.

Это подтверждает актуальность разработки веб-приложения с возможностью гибкой настройки и расширения функциональности.

\subsection{Анализ программных инструментов разработки}

Для реализации веб-приложения используется стек технологий, основанный на языке Python, что соответствует требованиям учебной дисциплины. Архитектура проекта включает серверную часть, клиентский интерфейс и систему шаблонизации.

Основные инструменты разработки:

\begin{itemize}
	\item \textbf{Backend-фреймворк:} Flask — лёгкий и гибкий веб-фреймворк на языке Python, предназначенный для разработки серверной части и REST API.
	
	\item \textbf{Язык программирования:} Python — используется для реализации бизнес-логики, обработки данных гардероба и формирования пользовательских образов.
	
	\item \textbf{Frontend:} HTML, CSS — применяются для построения пользовательского интерфейса и визуального представления данных. При необходимости используется JavaScript для динамического взаимодействия с интерфейсом.
	
	\item \textbf{Шаблонизация:} Jinja2 — используется для генерации HTML-страниц на стороне сервера и интеграции данных Python в пользовательский интерфейс.
	
	\item \textbf{База данных:} PostgreSQL — используется для хранения информации об элементах гардероба, пользователях и сформированных образах.
	
	\item \textbf{Работа с изображениями:} хранение фотографий одежды и образов осуществляется с использованием локального или облачного файлового хранилища.
\end{itemize}

Использование Flask в сочетании с Jinja2 позволяет реализовать серверно-рендеринг интерфейса, а также обеспечить простую интеграцию с HTML-шаблонами.

Применение HTML и CSS обеспечивает создание удобного и адаптивного пользовательского интерфейса, а Python позволяет эффективно реализовать обработку данных и алгоритмы формирования образов.

\subsection{Формулировка цели и задач работы}

На основе анализа существующих решений и выявленных ограничений сформулирована цель курсовой работы.

\textbf{Цель работы} --- разработка веб-приложения для управления личным гардеробом и формирования пользовательских образов с использованием фреймворка Flask, а также технологий HTML, CSS и Jinja2 для реализации пользовательского интерфейса.

Для достижения поставленной цели необходимо решить следующие задачи:

\begin{enumerate}
	\item Провести анализ существующих веб-сервисов управления гардеробом;
	\item Разработать архитектуру веб-приложения и структуру базы данных;
	\item Реализовать серверную часть приложения на Flask;
	\item Разработать пользовательский интерфейс с использованием HTML и CSS;
	\item Реализовать функциональность управления элементами гардероба (добавление, редактирование, удаление);
	\item Разработать механизм формирования пользовательских образов на основе одежды;
	\item Реализовать загрузку и хранение изображений;
	\item Провести тестирование разработанного веб-приложения.
\end{enumerate}